\section{Propósito y alcance del manual}

Este manual reúne prácticas y una formación inicial para que el equipo de AUTOVIDA incorpore accesibilidad y calidad durante el diseño, la implementación y la revisión del sistema. Está dirigido a todas las personas del equipo, no solo a quienes programan: las decisiones de contenido, diseño, pruebas y documentación también afectan a la experiencia de uso.

La primera versión desarrolla las pautas acordadas hasta ahora para la interfaz (colores, ausencia de ventanas emergentes, legibilidad del texto y compatibilidad con opciones de acceso como el barrido), además de criterios de accesibilidad para el backend y controles de calidad. Se ampliará conforme se definan nuevas necesidades con el centro y con las personas usuarias.

\subsection{Objetivos de aprendizaje}
Al terminar la formación, cada integrante debería poder:
\begin{itemize}[leftmargin=*]
  \item reconocer decisiones visuales y de interacción que pueden dificultar el uso;
  \item aplicar una paleta con contraste suficiente y no depender solo del color para comunicar;
  \item evitar interrupciones inesperadas, mantener los textos breves y comprobar que el barrido permite seleccionar correctamente los botones;
  \item permitir ajustar el tamaño del texto y de los botones, tanto para ampliarlos como para reducirlos, sin perder legibilidad ni facilidad de selección;
  \item proporcionar a cada botón un nombre accesible claro y añadir texto alternativo a los iconos cuando corresponda, para que las tecnologías de apoyo anuncien su función;
  \item aplicar criterios de accesibilidad también en el backend;
  \item revisar los cambios con criterios acordados, registrar incidencias y comprobar que se han resuelto.
\end{itemize}

\subsection{Principio de trabajo}
La accesibilidad y la calidad se incorporan desde el principio, no se reservan para una revisión final. Las pautas de este documento son una base de trabajo; no sustituyen las pruebas con alumnado y profesionales del centro ni la validación de requisitos concretos.

\begin{infobox}[title={Cómo utilizar este manual}]
Antes de implementar una funcionalidad, el equipo revisará las pautas aplicables y acordará cómo comprobarlas. En la revisión del cambio se anotarán las evidencias, los problemas detectados y quién se encarga de resolverlos.
\end{infobox}